\documentclass[11pt]{article}

\usepackage[margin=1in]{geometry}
\usepackage{amsmath, amssymb, amsthm}
\usepackage{enumitem}
\usepackage{hyperref}

\newtheorem{proposition}{Proposition}
\theoremstyle{definition}
\newtheorem{definition}{Definition}
\theoremstyle{remark}
\newtheorem*{remark}{Remark}

\newcommand{\Nin}{\mathcal{N}_i^{-}}
\newcommand{\J}{\mathcal{J}_i(t)}
\newcommand{\R}{\mathcal{R}}
\newcommand{\code}[1]{\texttt{#1}}

\title{The Shield Theorem\\[0.3em]\large A learned LF defender keeps the guarantees of Sundaram and Gharesifard}
\author{}
\date{}

\begin{document}
\maketitle

\begin{abstract}
A learned defender that only re-weights the values Local Filtering (LF) keeps, using a floored softmax, inherits every guarantee of Sundaram and Gharesifard~\cite{sg2019}. The result holds for \emph{every} parameter setting of the policy. It therefore also holds during training and against an attacker that has learned to exploit the policy.
\end{abstract}

\section{Setting}

We follow~\cite[Sec.~V]{sg2019}. Let $\R$ be the set of regular nodes and $\Nin$ the in-neighbours of node $v_i$, i.e.\ the nodes whose values $v_i$ receives. At each step $t$, every regular node $v_i$
\begin{enumerate}[label=(\roman*)]
  \item receives the values $\{x_j(t) : v_j \in \Nin\}$;
  \item removes the $F$ highest values above its own and the $F$ lowest values below its own (all of them if there are fewer than $F$), breaking ties arbitrarily; the retained in-neighbours form $\J \subseteq \Nin$;
  \item updates
  \begin{equation}\label{eq:lf}
    x_i(t+1) = a_{ii}(t)\,x_i(t) + \sum_{j \in \J} a_{ij}(t)\,x_j(t) - \alpha_t\, d_i(t),
  \end{equation}
  where $d_i(t)$ is a subgradient of $f_i$ at the convex combination $a_{ii}(t)\,x_i(t) + \sum_{j \in \J} a_{ij}(t)\,x_j(t)$.
\end{enumerate}
Equation~\eqref{eq:lf} is eq.~(8) of~\cite{sg2019}. The only requirement placed on the weights is that \emph{``at each time-step $t$ and for each $v_i \in \R$, the weights $a_{ij}(t)$, $v_j \in \{v_i\} \cup \J$, are lower-bounded by some strictly positive real number $\eta$ and sum to $1$''}~\cite[Sec.~V]{sg2019}. The weights may change over time and need not follow any formula; nothing in~\cite{sg2019} restricts how they are chosen.

\section{The shielded rule}

\begin{definition}[Shielded rule]\label{def:shield}
Fix $0 < c \le 1$. At each step $t$, node $v_i$ first computes $\J$ with the LF filter (step (ii) above). A policy $\pi_\theta$ then outputs finite scores $s_{ij}(t) \in \mathbb{R}$ for $j \in \{i\} \cup \J$. With
\[
  m_i(t) = |\J| + 1, \qquad \eta_i = \frac{c}{|\Nin| + 1},
\]
the weights are
\begin{equation}\label{eq:floored}
  a_{ij}(t) =
  \begin{cases}
    \eta_i + \bigl(1 - m_i(t)\,\eta_i\bigr)\,\dfrac{e^{s_{ij}(t)}}{\sum_{k \in \{i\}\cup\J} e^{s_{ik}(t)}}, & j \in \{i\} \cup \J,\\[1.2em]
    0, & j \in \Nin \setminus \J.
  \end{cases}
\end{equation}
The node then takes the subgradient step of~\eqref{eq:lf} at the resulting convex combination.
\end{definition}

The scores may depend on anything: the node's own state, the received values, its gradient, the step size, its history, or even the adversaries' values. Our implementation (\code{ShieldedDefender} in \code{advagent/learn.py}) uses the node's own state, the retained values, the degrees, the gradient and $\alpha_t$.

\section{Main result}

\begin{proposition}[Shield]\label{prop:shield}
Let $d_{\max} = \max_{i \in \R} |\Nin|$ and $\eta = c/(d_{\max}+1)$. For every $\theta$, and more generally for any rule that produces finite scores, the regular nodes running the shielded rule of Definition~\ref{def:shield} follow the LF dynamics~\eqref{eq:lf} with parameter $F$ and weight bound $\eta$. Consequently:
\begin{enumerate}[label=(\alph*)]
  \item \textbf{\cite[Thm.~6.1]{sg2019}.} If the adversaries form an $F$-local set of Byzantine nodes, the graph is $(2F+1)$-robust, the subgradients are bounded by $L$ and $\alpha_t \to 0$, then the regular nodes reach consensus.
  \item \textbf{\cite[Thm.~6.4]{sg2019}.} If the adversaries form an $F$-total set of malicious nodes and the graph is $(F+1,F+1)$-robust, then the regular nodes reach consensus.
  \item \textbf{\cite[Thm.~7.1]{sg2019}.} Under either of these conditions, if in addition $\sum_t \alpha_t = \infty$ and $\alpha_t \to 0$, every regular state converges to the interval
  \[
    \Bigl[\, \min_{i \in \R} \min \mathcal{M}_i,\ \max_{i \in \R} \max \mathcal{M}_i \,\Bigr],
  \]
  where $\mathcal{M}_i$ is the set of minimizers of $f_i$, whatever the adversaries do.
\end{enumerate}
All three conclusions hold for any learned policy, including one that is still training and one that an adversary has learned to exploit.
\end{proposition}

\begin{proof}
Fix an arbitrary step $t$ and regular node $v_i$, and write $p_{ij} = e^{s_{ij}(t)} / \sum_{k} e^{s_{ik}(t)}$ for the softmax term. Since the scores are finite, $p_{ij}$ is well defined, $p_{ij} > 0$ and $\sum_{j \in \{i\}\cup\J} p_{ij} = 1$. We check the four requirements of~\eqref{eq:lf}.

\emph{1. Support.} By~\eqref{eq:floored}, $a_{ij}(t)$ is nonzero only for $j \in \{i\} \cup \J$, and $\J$ is exactly the set that LF retains.

\emph{2. Lower bound.} Because $\J \subseteq \Nin$, we have $m_i(t) \le |\Nin| + 1$, so
\[
  1 - m_i(t)\,\eta_i \;\ge\; 1 - (|\Nin|+1)\,\eta_i \;=\; 1 - c \;\ge\; 0.
\]
Together with $p_{ij} \ge 0$ this gives
\[
  a_{ij}(t) \;\ge\; \eta_i \;=\; \frac{c}{|\Nin|+1} \;\ge\; \frac{c}{d_{\max}+1} \;=\; \eta \;>\; 0
  \qquad \text{for all } j \in \{i\}\cup\J.
\]

\emph{3. Sum to one.}
\[
  \sum_{j \in \{i\}\cup\J} a_{ij}(t) \;=\; m_i(t)\,\eta_i + \bigl(1 - m_i(t)\,\eta_i\bigr) \sum_{j \in \{i\}\cup\J} p_{ij} \;=\; 1.
\]

\emph{4. Subgradient point.} The subgradient $d_i(t)$ is evaluated at the convex combination formed with these weights, as~\eqref{eq:lf} requires (\code{torchsim.step} evaluates the gradient there).

Hence the closed loop is an instance of~\eqref{eq:lf} with parameter $F$ and bound $\eta$.

It remains to see that the theorems apply. The hypotheses of~\cite[Thms.~6.1, 6.4, 7.1]{sg2019} concern only the graph, the adversary model, the functions $f_i$, the step sizes, and the fact that the regular nodes run~\eqref{eq:lf} with weights bounded below by $\eta$. Their proofs use only per-step properties of the weights: the reduction to the regular-only dynamics (9) in~\cite[Prop.~5.1]{sg2019} uses $a_{ij}(t) \ge \eta$ and $\sum_j a_{ij}(t) = 1$, and the consensus arguments use lower bounds of $\eta$ or $\eta/2$ on the entries of the equivalent matrix $\bar{A}(t)$. None of them uses how the weights were chosen, so the conclusions (a)--(c) hold.
\end{proof}

\section{Remarks}

\begin{remark}[What is not inherited]
\cite{sg2019} gives no optimality guarantee for LF, and none is claimed here. The lower bound of~\cite[Prop.~8.1]{sg2019} applies to the shielded rule as it does to any rule that keeps the regular states in the hull: on some instances the error is at least $\frac{|T|}{n}\,|b-a|$, where $T$ is a maximum $F$-local set. The learned weights can only move the result \emph{within} the guaranteed set.
\end{remark}

\begin{remark}[Relation to plain LF]
If all scores of node $v_i$ are equal, then $p_{ij} = 1/m_i(t)$ and
\[
  a_{ij}(t) = \eta_i + \frac{1 - m_i(t)\,\eta_i}{m_i(t)} = \frac{1}{m_i(t)},
\]
which is exactly LF's uniform weight over $\{i\}\cup\J$, for any $c$. Our policy's last layer is initialised to zero, so training starts from plain LF.
\end{remark}

\begin{remark}[The trade-off set by $c$]
The policy controls only the mass $1 - m_i(t)\,\eta_i$ above the floor. A larger $c$ raises the floor $\eta$, which strengthens the consensus rate in the proofs but leaves the policy less freedom; when LF removes nothing ($m_i = |\Nin|+1$), $c = 1$ leaves no freedom at all. A smaller $c$ gives the policy more freedom at the price of a smaller $\eta$. We use $c = 0.25$, so every weight is at least a quarter of $1/(|\Nin|+1)$.
\end{remark}

\begin{remark}[Implementation guard]
The proof needs finite scores, since non-finite scores would make the softmax undefined. The implementation therefore replaces NaN and $\pm\infty$ scores before the softmax (\code{torch.nan\_to\_num}).
\end{remark}

\begin{remark}[Unshielded ablation]
\code{ShieldedDefender(shield=False)} drops the filter (step (ii)) and the floor. It is not an instance of~\eqref{eq:lf}, so Proposition~\ref{prop:shield} does not apply. In our experiments a constant $+50$ attack captures it (mean error $39$ for $F=1$), while the shielded defender is never unsafe.
\end{remark}

\begin{remark}[Empirical check]
\code{tests/test\_phase5.py} runs randomly initialised policies with large parameters, far from LF, against the strongest scripted attacks and checks hull containment and consensus. The Phase~5--7 evaluations report zero safety violations for every trained shielded policy.
\end{remark}

\begin{thebibliography}{9}
\bibitem{sg2019}
S.~Sundaram and B.~Gharesifard, ``Distributed optimization under adversarial nodes,'' \emph{IEEE Trans. Autom. Control}, vol.~64, no.~3, pp.~1063--1076, Mar.~2019.
\end{thebibliography}

\end{document}
